\section{Related Works}
Wornow et al.~\cite{wornow2023} introduced EHRSHOT, a longitudinal benchmark with tasks for evaluating clinical foundation models. We reuse its structured records and canonical patient split and define a new endpoint from the first recorded diagnosis in a stable community core. The task tests whether network features add value for community entry beyond diagnosis histories.

Crisafulli et al.~\cite{crisafulli2026} statistically validated disease co-occurrence edges and compared communities across demographic groups. Our Project~1 graph likewise filters co-occurrences and supplies the communities and bridge features for the prediction task. Singh et al.~\cite{singh2026} used structural and demographic features in GCN, GraphSAGE, and GAT models to predict obstetric--non-obstetric co-occurrence links. We apply the same architecture families to patient-landmark classification, where the outcome is a later recording and the graph edges describe population-level co-occurrence.

Chet learns diagnosis representations on dynamic disease graphs that follow visit-to-visit changes~\cite{lu2022chet}. GraphCare augments individual patient graphs with external biomedical knowledge for tasks such as readmission and mortality prediction~\cite{jiang2024graphcare}. Both inform our patient-level use of graph context. Our graph comes from observed co-occurrences, and controlled edge rewiring tests the added value of that topology while retaining node features.

Oss Boll et al.~\cite{boll2024} found diagnosis prediction to be the most common task in their survey of EHR GNNs, with MIMIC-III the most common data source. Our comparison extends that setting to longitudinal EHRSHOT records and diagnosis groups defined by stable graph communities. Hancox et al.~\cite{hancox2024} documented limited calibration reporting and frequent risk-of-bias concerns in graph-based clinical prediction. We address those evaluation gaps with patient-disjoint splits, matched baselines, and calibration alongside discrimination for each target.
